% === I. INTRODUCTION =========================================================
\section{Introduction}

\IEEEPARstart{M}{odel-based} controllers for autonomous racing derive their
performance from the fidelity of the vehicle model they carry
\cite{liniger2015,kabzan2019lbmpc,forzaeth2024}. In the lateral channel that
fidelity is dominated by the tire, simultaneously the physical bound on
manoeuvrability and the least accessible component to measure
\cite{pacejka2012book,amz2020}. Classical characterisation takes the vehicle
off the track for steady-state circular or ramp-steer manoeuvres, at a
logistical cost a race weekend rarely affords; identifying on track removes the
logistics but brings back transients, noise and a racing line that is a poor
excitation signal, and nonlinear least squares applied directly to lap data is
brittle in exactly this regime \cite{ontrack2025}.

Dikici \emph{et al.}~\cite{ontrack2025} propose the hybrid remedy that is the
starting point of this paper: a small residual network trained on
\SI{30}{\second} of ordinary lap data predicts the one-step error of a nominal
single-track model; the corrected model is rolled out through a \emph{synthetic}
quasi-steady steering sweep; Pacejka coefficients are fitted to that sweep; and
each fit becomes the next nominal model. Because the sweep is generated rather
than driven it can be made steady and noise-free, and the residual network only
has to interpolate within its training distribution. On a 1:10 platform the
procedure reports a \num{3.3}$\times$ lower one-step RMSE than nonlinear least
squares under noise \cite{ontrack2025}.

\subsection{Problem statement}

This paper asks what the pipeline identifies when it is moved from a
\SI{3.7}{\kilo\gram}, \SI{0.32}{\meter}-wheelbase platform to a full-scale
Formula Student vehicle: \SI{269.6}{\kilo\gram}, \SI{1.53}{\meter} wheelbase,
\SI{16}{\degree} of road-wheel lock, following a reference speed profile of
\SIrange{9.6}{25.8}{\meter\per\second}. The question is identifiability, not
parameter transfer, and two scale-dependent mechanisms, both latent at 1:10,
break the estimate. First, excitation collapses inside the virtual-data
step: the Pacejka fit never sees the recorded lap, only the synthetic rollout,
whose reachable friction
$\mu_{\mathrm{reach}}\!\approx\!v^{2}\delta_{\mathrm{sweep}}/(gL)$ is a property
of the rollout configuration and not of the tire, and a rollout speed inherited
from the scaled platform caps it at $0.43$ here. Below the tire's peak the Magic
Formula degenerates to $F_{y}\approx F_{z}(BCD)\alpha$, only the product $BCD$
is identifiable, and the bounded optimiser resolves the remaining freedom on a
box edge \cite{ljung1999sysid,pacejka2012book}. Second, lateral demand
scales as $v^{2}$, so the tracking cost of a bad grip estimate grows
quadratically with speed:
one identified set with a peak friction coefficient of \num{0.40} tracked at
\SI{0.16}{\meter} of cross-track error at \SI{11.1}{\meter\per\second} and
\SI{41.89}{\meter} at \SI{27}{\meter\per\second} through the same controller
(Section~\ref{subsec:gates}), and the scaled platform never reaches the speed at
which that error becomes observable.

One difference works in our favour. Because the study runs in the CARLA
simulator \cite{carla2017} through a purpose-built ROS~2 interface, the
simulator's own per-wheel telemetry (slip angle, normal load and lateral force)
is available as ground truth, so an on-track identifier can be scored against
the forces the plant actually generated.

\subsection{Contributions}

\begin{enumerate}
\item \textbf{An identifiability analysis of the virtual-data step}
(Section~\ref{subsec:excitation}): the closed-form reachable-friction bound
above, and an \emph{excitation-aware rollout} that selects the rollout speed
from the measured lap and warns when the configuration cannot reach the prior's
peak. A grid over rollout speed, sweep amplitude and true peak factor on a known
tire verifies it (Section~\ref{subsec:sweep}). This corrects the baseline
method~\cite{ontrack2025} at any scale, and stays latent at 1:10 because a
\SIrange{2}{4}{\meter\per\second} rollout is representative there and is not
here.

\item \textbf{Two structural properties of the loop, and the corrections they
imply} (Section~\ref{subsec:contraction}): stage~3's residual depends on $B$,
$C$ and $D$ only through their product, so the peak factor is not recoverable
from the rollout \emph{at any excitation} and must be measured on the recorded
trace and held; and the loop is a positive feedback path rather than a
contraction, so stage~4 becomes a relaxed fixed-point iteration at
$\beta=\num{0.20}$.

\item \textbf{The corrections the bound implies}, on both sides of the fit
(Sections~\ref{subsec:iterative} and~\ref{subsec:measurement}): a sweep bounded
at $1.5\times$ the 99th percentile of the observed steering, a Tikhonov target
frozen at the configured prior, slip angles built from the \emph{achieved}
road-wheel angle, mass and wheelbase from measured static wheel loads, a
friction warm start from the \emph{curvature} of the axle force curve, and
physics-informed dynamics residuals normalised to the state-increment units of
the data term, which the published form leaves a factor \num{2.3e5} apart.

\item \textbf{Admissibility gates on derived physical quantities} at the
model-to-controller boundary (Section~\ref{sec:integration}): axle cornering
stiffness, understeer sign and critical speed, and grip ceiling against the
trajectory's lateral demand. They replace per-coefficient box checks, which
cannot detect a degenerate fit or a pairwise-oversteering axle set.

\item \textbf{A simulation architecture that supplies the ground truth, and a
closed-loop evaluation on it} (Sections~\ref{subsec:platform}--\ref{subsec:data}
and~\ref{sec:results}). The evaluation runs on a static \num{0.70} friction
surface, deliberately not the \num{1.05} the configured tire prior was measured
on, so a pipeline unable to move off its prior scores badly. Over three
timed laps the corrected pipeline identifies axle cornering stiffness to
\SI{3.2}{\percent} front and \SI{1.4}{\percent} rear and tracks the peak factor
to \SI{11}{\percent}, against an inherited configuration that rails $D$ on its
\num{2.0} bound in every cycle and overstates stiffness by
\SIrange{118}{132}{\percent}. Axle force RMSE is \num{12.1}$\times$ and
\num{10.5}$\times$ lower, and MPC lateral tracking error \SI{32}{\percent}
lower, at unchanged solver cost.
\end{enumerate}

All results are obtained in simulation, and Section~\ref{sec:disc} states which
conclusions are expected to transfer to hardware.
